% Created 2026-05-12 ter 18:03
% Intended LaTeX compiler: lualatex
\documentclass[12pt,a4paper]{article}
\usepackage{fgv-paper}
% Bibliografia controlada pelo TOML; os estilos FGV não carregam biblatex.
\usepackage[backend=biber,sorting=nty,giveninits=true]{biblatex}

\addbibresource{resumo_artigos_encontro_3.bib}
\geometry{a4paper,left=3cm,right=2cm,top=3cm,bottom=2cm,headheight=58pt,headsep=18pt,footskip=1.2cm}
\usepackage{float}
\usepackage{tikz}
\usepackage{tcolorbox}
\usepackage{tabularray}
\usepackage{ltablex}
\keepXColumns
\setlength{\emergencystretch}{3em}
\sloppy
\definecolor{FGVHeaderBlueDark}{HTML}{003B71}
\definecolor{FGVHeaderBlueMid}{HTML}{006BB6}
\definecolor{FGVHeaderBlueLight}{HTML}{8FD3F4}
\setlength{\headheight}{58pt}
\setlength{\headsep}{18pt}
\newcommand{\FGVHeaderLogoInclude}{\IfFileExists{/home/gustavodetarso/Documentos/mppg/software/academic_pipeline_mppg/app_bundle/institutions/fgv/assets/fgv.png}{
\includegraphics[height=1.05cm]{/home/gustavodetarso/Documentos/mppg/software/academic_pipeline_mppg/app_bundle/institutions/fgv/assets/fgv.png}}{}}
\newcommand{\FGVHeaderGradientRule}{\begin{tikzpicture}[baseline=(current bounding box.center)]\shade[left color=FGVHeaderBlueDark,middle color=FGVHeaderBlueMid,right color=FGVHeaderBlueLight] (0pt,0pt) rectangle (\headwidth,2pt);\end{tikzpicture}}
\newcommand{\FGVHeaderBlock}{\makebox[\headwidth][l]{\begin{minipage}[t]{\headwidth}\FGVHeaderLogoInclude\\[-1pt]\FGVHeaderGradientRule\end{minipage}}}
\fancypagestyle{fgvreportstyle}{%
\fancyhf{}%
\fancyhead[L]{\FGVHeaderBlock}%
\renewcommand{\headrulewidth}{0pt}%
\fancyfoot[C]{\thepage}%
}
\fancypagestyle{plain}{%
\fancyhf{}%
\fancyhead[L]{\FGVHeaderBlock}%
\renewcommand{\headrulewidth}{0pt}%
\fancyfoot[C]{\thepage}%
}
\AtBeginDocument{\pagestyle{fgvreportstyle}\thispagestyle{fgvreportstyle}}
\institution{fgv}
\programname{}
\coursename{Mestrado Acadêmico em Políticas Públicas e Governo}
\disciplinename{Políticas Públicas - Teoria e Análise}
\professorname{Lizandro Lui}
\cityname{Brasília}
\author{Gustavo M. Mendes de Tarso}
\date{\today}
\title{Resumo analítico dos textos}
\hypersetup{
 pdfauthor={Gustavo M. Mendes de Tarso},
 pdftitle={Resumo analítico dos textos},
 pdfkeywords={},
 pdfsubject={},
 pdfcreator={Emacs 28.2 (Org mode 9.5.5)}, 
 pdflang={Pt_Br}}
\begin{document}

\begingroup
\linespread{1}\selectfont
\begin{tcolorbox}[title=Ficha Técnica,
  colback=gray!5,colframe=gray!40,boxrule=0.4pt,sharp corners]
\begin{tblr}{rowsep=1pt,stretch=1, rows={t}, colspec={Q[l,2.8cm] X[l]}}
\textbf{Curso}                   & Mestrado Acadêmico em Políticas Públicas e Governo \\
\textbf{Turma}                   & 2026.1 \\
\textbf{Pólo}                    & Brasília \\
\textbf{Disciplina}              & Políticas Públicas - Teoria e Análise \\
\textbf{Professor}               & Lizandro Lui \\
\textbf{Aluno(s)}                & Gustavo M. Mendes de Tarso \\
\textbf{Data}                    & 2026 \\
\textbf{Título do trabalho}      & Resumo analítico dos textos \\
\end{tblr}
\end{tcolorbox}
\endgroup

\clearpage

\section{Introdução}
\label{sec:org7d55e60}

O corpus reúne quatro textos que, embora não formem um conjunto tematicamente homogêneo em sentido estrito, convergem em torno de um problema disciplinar comum: compreender a implementação de políticas públicas como processo efetivo de produção da ação estatal. Em vez de tratar a implementação como etapa residual entre formulação e resultados, os trabalhos a examinam em planos distintos e complementares: a coordenação federativa e os efeitos de sistemas nacionais sobre a provisão local \parencite{bichir2020_sistemas_nacionais}; a consolidação teórica e analítica do próprio campo dos estudos de implementação no Brasil \parencite{lotta2019_teorias_analises}; as práticas cotidianas de categorização de usuários por agentes implementadores \parencite{lotta2020_categorizando_usuarios}; e as interfaces entre Estado e organizações confessionais na prestação de serviços públicos \parencite{lotta2023_jesus_speaks}.

O foco de leitura adotado neste documento é, portanto, menos substantivo-setorial do que analítico. O interesse principal não está em extrair uma unidade temática artificial entre assistência social, saúde, coordenação federativa e cooperação com organizações religiosas, mas em identificar como cada texto ilumina dimensões específicas da implementação: suas mediações institucionais, suas escalas de operação, seus mecanismos de coordenação, seus julgamentos práticos e seus efeitos sobre a distribuição concreta de bens e serviços públicos. Essa opção é coerente com a própria heterogeneidade do corpus, que combina obra de enquadramento teórico, estudo de alcance nacional, investigação microinteracional e análise interorganizacional.

A unidade possível entre os textos não reside, assim, em um mesmo objeto empírico, mas em um diagnóstico comum: políticas públicas não se realizam automaticamente por força do desenho formal. Entre norma e prática, há burocracias, capacidades, arranjos federativos, valores, classificações e interfaces organizacionais que transformam o conteúdo concreto da ação estatal. A coletânea organizada por Gabriela Lotta explicita esse deslocamento ao afirmar a implementação como campo autônomo de investigação \parencite{lotta2019_teorias_analises}; os demais textos mostram, em registros diferentes, como essa autonomia analítica se confirma empiricamente, seja no efeito estruturante de sistemas nacionais \parencite{bichir2020_sistemas_nacionais}, na produção cotidiana de diferenças sociais \parencite{lotta2020_categorizando_usuarios} ou na construção de ajustes entre lógicas institucionais heterogêneas \parencite{lotta2023_jesus_speaks}.

O objetivo deste documento é oferecer uma leitura de fechamento capaz de articular o corpus em chave interpretativa, consolidando seus eixos transversais sem reduzir suas diferenças. Isso implica, de um lado, reconstruir a lógica comum que atravessa os textos; de outro, preservar a especificidade de cada contribuição quanto ao problema abordado, à escala de análise, ao método mobilizado e ao tipo de mecanismo explicativo privilegiado. A aposta analítica é que o conjunto permite compreender a implementação como fenômeno multiescalar: ela é simultaneamente estruturada por instituições, operada por organizações, mediada por burocracias e experimentada nas relações concretas entre Estado e públicos.

O caminho analítico adotado acompanha esse diagnóstico. Primeiro, parte-se do reconhecimento da heterogeneidade do corpus para evitar comparações forçadas. Em seguida, privilegiam-se apenas os eixos de comparação efetivamente sustentados pelos textos: problema de pesquisa, arquitetura conceitual, escala analítica, tipo de evidência, mecanismos explicativos, contribuição para a disciplina e limites. Por fim, a síntese analítica desloca a leitura da comparação descritiva para uma interpretação própria: a de que o principal ganho do conjunto está em demonstrar que a implementação é o lugar em que a política pública deixa de ser apenas desenho e se torna relação social institucionalizada, com efeitos distributivos, organizacionais e normativos que variam conforme os arranjos pelos quais o Estado opera.



\section{Resumo individual dos textos}
\label{sec:org42c6cc6}



\subsection{Sistemas nacionais de políticas públicas e seus efeitos na implementação: o caso do Sistema Único de Assistência Social (Suas)}
\label{sec:org850e4f1}

O artigo toma como problema central uma questão clássica do federalismo brasileiro, mas a desloca para um terreno empírico mais preciso: não basta afirmar genericamente que a União coordena e os entes subnacionais implementam; é necessário investigar como arranjos institucionais específicos alteram, de fato, a provisão local de políticas sociais em um país marcado por heterogeneidade territorial, desigualdade de capacidades administrativas e autonomia municipal. A pergunta que organiza o texto é, portanto, se um sistema nacional de política pública, no caso o Suas, produz efeitos observáveis sobre a implementação municipal dos serviços socioassistenciais e sobre a redução de desigualdades entre municípios. Em vez de permanecer no plano normativo da engenharia federativa, o artigo procura mostrar como regras, incentivos e instrumentos de coordenação se traduzem — ou não — em convergência efetiva na ponta da política pública\parencite{bichir2020_sistemas_nacionais}

O objetivo e o escopo do estudo são delimitados com clareza. Os autores não pretendem reconstituir toda a trajetória de institucionalização da assistência social nem oferecer uma história geral do Suas. O foco recai sobre os efeitos das normas e instrumentos federais na implementação dos serviços socioassistenciais básicos entre 2011 e 2016, com atenção especial a duas dimensões: o incremento da oferta e a convergência territorial de padrões mínimos de provisão. Em complemento, o artigo busca aferir se maiores níveis de implementação estão associados a capacidades institucionais locais, operacionalizadas por meio de aspectos de recursos humanos e infraestrutura dos Centros de Referência de Assistência Social. Esse recorte é decisivo porque permite examinar simultaneamente a força da indução federal e a persistência de variações subnacionais, evitando tanto a leitura de que a descentralização dissolveria qualquer coordenação quanto a visão oposta de que a centralização regulatória explicaria sozinha os resultados observados\parencite{bichir2020_sistemas_nacionais}

A tese principal do artigo pode ser sintetizada em dois movimentos articulados. O primeiro é que sistemas nacionais de políticas públicas, quando combinam regulação central, transferências condicionadas, financiamento fundo a fundo, instâncias de pactuação federativa e definição de parâmetros nacionais, possuem capacidade de induzir a expansão da implementação local e reduzir desigualdades territoriais. O segundo é que esse efeito não elimina a relevância das capacidades locais: ao contrário do que sugere parte da literatura que enfatiza predominantemente a indução federal, os autores mostram que a implementação dos serviços básicos do Suas se associa positivamente, ainda que de forma fraca, à estruturação administrativa dos equipamentos locais. A implicação analítica dessa formulação é importante: a coordenação nacional não apaga a política local, mas a reorganiza dentro de parâmetros mais densos de regulação e incentivo. O resultado é um arranjo de governança multinível em que a União é forte para definir direção, mas os municípios continuam relevantes para converter normas em provisão concreta\parencite{bichir2020_sistemas_nacionais}

No plano conceitual, o texto mobiliza categorias centrais para a literatura de políticas públicas no Brasil. A mais importante é a de sistemas nacionais de políticas públicas, entendidos como arranjos institucionais que estruturam a descentralização sob coordenação federal, definindo agendas, padrões e incentivos para implementação subnacional. O artigo também trabalha com as noções de coordenação federativa, indução federal, descentralização nacionalmente regulada, capacidades institucionais e governança multinível. Essas categorias não aparecem como mera moldura terminológica: elas organizam a interpretação dos dados e servem para qualificar a assistência social como um caso particularmente exigente. Diferentemente de políticas mais facilmente padronizáveis, os serviços socioassistenciais dependem intensamente de trabalho relacional, presença territorial, infraestrutura física e equipes profissionais, o que torna a implementação mais sensível às condições locais. Daí a relevância de testar empiricamente até que ponto a arquitetura nacional do Suas consegue produzir convergência em um domínio que, por natureza, resiste à uniformização simples\parencite{bichir2020_sistemas_nacionais}

A reconstrução do objeto feita pelos autores é uma das forças do artigo. O Suas é apresentado não apenas como desenho institucional formal, mas como resposta a um legado histórico de fragmentação, filantropia, benemerência e baixa rotinização estatal da assistência social. Esse ponto é substantivo porque esclarece por que a simples existência de normas federais não poderia produzir efeitos imediatos nem homogêneos. O texto insiste, corretamente, em uma temporalidade institucional: regras produzem efeitos ao longo do tempo, acumulam capacidades, sedimentam expectativas e redefinem incentivos. Assim, o caso do Suas é analiticamente fecundo porque permite observar um processo de construção relativamente recente de estatalidade em uma área antes marcada por provisão dispersa e pouco regulada. A ênfase no tempo institucional impede leituras instantaneístas da implementação e reforça a ideia de que a avaliação dos efeitos do sistema nacional deve considerar trajetórias de institucionalização, não apenas fotografias normativas\parencite{bichir2020_sistemas_nacionais}

Do ponto de vista metodológico, o artigo se diferencia por enfrentar uma lacuna relevante em um campo frequentemente dominado por estudos de caso. Em lugar de uma análise local ou exclusivamente qualitativa, os autores constroem indicadores comparáveis para examinar padrões nacionais de implementação e de capacidades estatais no período de 2011 a 2016. A principal estratégia consiste em adaptar informações do Censo Suas e do Indicador de Desenvolvimento do Centro de Referência de Assistência Social para compor medidas capazes de capturar, ao longo do tempo, a oferta de serviços socioassistenciais básicos e o grau de estruturação das unidades responsáveis por essa oferta. O procedimento não é apenas técnico; ele expressa uma escolha teórica: medir implementação por processos e pela presença de estruturas de provisão, e não apenas por resultados finais, o que é particularmente adequado em uma política cuja mensuração de resultados é mais recente e menos consolidada. O artigo, assim, avança ao transformar um debate frequentemente abstrato sobre coordenação federativa em uma investigação comparativa baseada em evidências sistemáticas\parencite{bichir2020_sistemas_nacionais}

Os achados apresentados sustentam a hipótese de que houve incremento da implementação dos serviços básicos do Suas e, ao mesmo tempo, maior convergência entre municípios, regiões e portes populacionais. Isso significa que o sistema nacional não apenas expandiu a presença da política, mas também contribuiu para reduzir a distância relativa entre contextos municipais distintos. O argumento é particularmente relevante para a disciplina porque oferece evidência a favor da ideia de que a centralização regulatória pode ser compatível com redução de desigualdades territoriais em contextos federativos. Ao mesmo tempo, os autores não transformam essa conclusão em celebração simplista da União. A associação positiva entre implementação e capacidades locais, ainda que modesta, mostra que o efeito do sistema depende de condições de recepção, organização administrativa e materialização institucional nos equipamentos municipais. Em termos substantivos, o artigo sugere que a indução federal cria as condições de possibilidade da expansão, mas não substitui o trabalho organizacional necessário para sustentar a política na ponta\parencite{bichir2020_sistemas_nacionais}

A principal contribuição do texto para os estudos de implementação está em reposicionar o debate para além da oposição entre formulação central e execução local. O artigo demonstra que a implementação deve ser pensada como efeito de arranjos institucionais multiníveis, nos quais normas federais, incentivos financeiros, padrões administrativos e capacidades municipais interagem. Nesse sentido, ele ajuda a sofisticar a compreensão de implementação ao aproximá-la de questões de coordenação federativa e de construção incremental de capacidades estatais. Para a administração pública e para a análise de políticas sociais, isso tem duas implicações importantes. Primeiro, mostra que desenhos institucionais importam e podem alterar padrões territoriais de provisão. Segundo, indica que a qualidade da implementação não é redutível à adesão formal ao sistema: depende de recursos humanos, infraestrutura e arranjos organizacionais concretos. A contribuição do artigo, portanto, é ao mesmo tempo teórica e aplicada, pois oferece uma moldura explicativa para pensar como políticas descentralizadas podem ser nacionalmente estruturadas sem eliminar variação local\parencite{bichir2020_sistemas_nacionais}

Os limites e tensões do estudo decorrem, em parte, de sua própria virtude analítica. Ao privilegiar indicadores de processo e estrutura para captar implementação, o artigo ilumina a expansão e a convergência da oferta, mas deixa em segundo plano a qualidade substantiva das interações entre burocracias e usuários, bem como os efeitos distributivos e simbólicos da provisão. Em outras palavras, saber que serviços se expandiram e que capacidades mínimas se disseminaram não resolve, por si, a questão de como os cidadãos experienciam essa política, nem se a padronização institucional reduz discricionariedades cotidianas ou, ao contrário, convive com novas formas de diferenciação no atendimento. Além disso, a correlação fraca entre implementação e capacidades locais é interpretada com cautela pelos próprios autores, o que abre espaço para indagar fatores não incorporados ao modelo, como trajetórias políticas locais, estilos de gestão, relações com organizações da sociedade civil e dinâmicas territoriais específicas. O artigo reconhece, assim, que “nem tudo começa e termina” no centro federal, mas também não esgota a explicação dos mecanismos subnacionais que mediam os efeitos do sistema nacional\parencite{bichir2020_sistemas_nacionais}

Há ainda uma tensão produtiva entre o argumento da convergência institucional e a natureza da própria assistência social. Como os autores sugerem, trata-se de uma política composta por serviços de baixa padronização e alta dependência de trabalho humano situado. Isso significa que a convergência medida por indicadores administrativos pode coexistir com forte variação nas práticas concretas de implementação. Essa observação não invalida os resultados; ao contrário, indica a fronteira de sua generalização. O artigo é muito convincente ao demonstrar efeitos sistêmicos sobre expansão e redução de desigualdades em termos de provisão e estrutura. Mas ele deixa aberta a agenda para investigar se essa convergência institucional se traduz em experiências similares de acesso, acolhimento, classificação e tratamento dos usuários nos territórios. A passagem dos macroefeitos do sistema para a micropolítica da implementação permanece, portanto, como agenda posterior, especialmente relevante em políticas baseadas em interação cotidiana\parencite{bichir2020_sistemas_nacionais}

Nesse ponto, o diálogo com o restante do corpus é especialmente fértil, desde que não se force uma unidade inexistente. Enquanto este artigo opera em escala sistêmica e federativa, examinando como arranjos nacionais moldam padrões de implementação municipal, outros textos do conjunto deslocam o olhar para níveis analíticos distintos. A abordagem sobre implementação reunida em coletânea organizada por Gabriela Lotta amplia precisamente esse repertório ao tratar a implementação como campo teórico plural, atento a burocracias, interações, discricionariedade, capacidades e contextos institucionais, o que ajuda a situar o artigo em uma agenda mais ampla de análise\parencite{lotta2019_teorias_analises,bichir2020_sistemas_nacionais}

O contraste mais nítido, porém, aparece quando se compara esta análise dos efeitos de um sistema nacional com estudos que enfatizam a produção cotidiana de diferenças na ponta. Ao mostrar que a implementação pode convergir institucionalmente sem eliminar mediações locais, o artigo cria uma ponte importante com trabalhos que investigam como burocratas categorizam usuários e produzem tratamentos diferenciados nas interações concretas do atendimento. A relação entre esses textos não é de contradição, mas de escalas analíticas complementares: aqui se demonstra a capacidade do arranjo nacional de estruturar a oferta; ali se revelam os mecanismos pelos quais essa oferta é filtrada, interpretada e distribuída no cotidiano. De modo semelhante, análises sobre interações entre Estado e organizações de base religiosa em contextos de implementação lembram que a provisão pública pode envolver hibridizações institucionais e ajustes situados que escapam a uma leitura puramente estatal do sistema. Assim, o artigo de Renata Bichir, Sergio Simoni Junior e Guilherme Pereira oferece uma peça crucial do quebra-cabeça do corpus: ele explica como a arquitetura federativa produz condições de implementação em escala nacional, mas o próprio corpus sugere que compreender integralmente a política exige descer também ao nível das interações, classificações e acomodações organizacionais que dão conteúdo efetivo a essa arquitetura\parencite{lotta2020_categorizando_usuarios,lotta2023_jesus_speaks,bichir2020_sistemas_nacionais}



\subsection{Teorias e Análises sobre Implementação de Políticas Públicas no Brasil}
\label{sec:org0c63bb0}

A obra organizada por Gabriela Lotta ocupa uma posição particular no corpus porque não se limita a examinar um caso substantivo de implementação, mas procura consolidar e dar inteligibilidade a um campo de pesquisa em expansão no Brasil. Seu problema central não é apenas descrever a implementação como etapa do ciclo de políticas, e sim enfrentar a tendência, recorrente na literatura e no debate aplicado, de tratar “implementação” como simples execução administrativa. Contra esse rebaixamento analítico, o livro sustenta que a implementação constitui um objeto próprio de investigação, no qual decisões, interações, mediações institucionais, capacidades estatais, discricionariedade e efeitos não previstos se tornam visíveis. O ponto de partida é que a política pública “como ela é” não coincide integralmente com a política desenhada, e que compreender essa distância exige categorias específicas de análise. \parencite{lotta2019_teorias_analises}

O objetivo e o escopo do livro são duplos. Em um primeiro nível, a coletânea funciona como mapa introdutório e agenda de pesquisa: apresenta recortes teóricos, metodológicos e temáticos dos estudos de implementação, situando-os no interior do campo mais amplo de análise de políticas públicas. Em um segundo nível, o volume busca demonstrar empiricamente a fecundidade dessa agenda a partir de capítulos que tratam de educação, assistência social, coordenação intergovernamental, avaliação, capacidades estatais, burocracias e consequências imprevistas. O livro, portanto, não defende uma teoria única da implementação; seu escopo é mais plural e programático. Ele organiza um espaço de interlocução entre abordagens que compartilham a recusa de tomar a política formulada como equivalente à política efetivamente produzida na prática. \parencite{lotta2019_teorias_analises}

A tese principal que emerge do prefácio e, sobretudo, do capítulo de abertura é que os estudos de implementação ampliam a capacidade explicativa da análise de políticas públicas ao deslocarem o foco do desenho formal para os processos concretos de produção da ação estatal. Em vez de conceber a implementação como fase residual, posterior e subordinada à formulação, o livro a trata como momento constitutivo da própria política. Isso implica reconhecer que objetivos, instrumentos e diretrizes são retrabalhados ao longo da execução; que atores implementadores não apenas obedecem, mas interpretam, priorizam e ajustam; e que resultados decorrem de arranjos, interações e contextos institucionais. A implementação, assim, não aparece como desvio acidental, mas como dimensão ordinária da vida das políticas públicas. \parencite{lotta2019_teorias_analises}

Do ponto de vista conceitual, a obra mobiliza um conjunto de categorias que estruturam o campo. Em primeiro lugar, a própria noção de ciclo de políticas é empregada como instrumento analítico, mas sem ingenuidade descritiva: trata-se de um recurso heurístico para distinguir agenda, formulação, implementação e avaliação, não de uma sequência linear que espelharia a realidade. Em segundo lugar, a implementação é definida por sua atenção às diferenças entre o formulado e o executado, às mediações institucionais que produzem essas diferenças e aos atores encarregados da materialização da política. Em terceiro lugar, ganham centralidade categorias como burocratas de nível de rua, burocracias de médio escalão, discricionariedade, arranjos institucionais, coordenação e capacidades estatais. O livro sugere que essas categorias não são acessórios temáticos, mas lentes para entender como o Estado opera concretamente. \parencite{lotta2019_teorias_analises}

Entre essas categorias, a figura dos burocratas de nível de rua e o conceito de discricionariedade ocupam lugar decisivo. O capítulo introdutório enfatiza que a ação estatal é produzida em contextos de incerteza, escassez, heterogeneidade de demandas e necessidade de julgamento situado. Nesses contextos, implementadores não apenas aplicam regras, mas as interpretam em interação com usuários, territórios e organizações. A discricionariedade, portanto, não é tratada apenas como problema de controle, mas como condição constitutiva da implementação. Essa formulação é importante porque desloca a análise da dicotomia entre fidelidade e infidelidade ao desenho original: mais relevante do que medir aderência normativa é compreender os mecanismos pelos quais escolhas situadas produzem inclusão, exclusão, adaptação e efeitos distributivos concretos. \parencite{lotta2019_teorias_analises}

O método ou, mais precisamente, a estratégia argumentativa da coletânea é coerente com esse propósito de consolidação de campo. O capítulo de abertura realiza uma síntese da literatura e apresenta seus pressupostos contemporâneos, enquanto os demais capítulos funcionam como demonstrações analíticas aplicadas. Em vez de um desenho único de pesquisa, o livro aposta na combinação de reconstrução teórica com pluralidade empírica: capítulos conceituais convivem com surveys, estudos de caso, análises de arranjos institucionais e investigações sobre políticas setoriais específicas. Essa heterogeneidade não constitui fraqueza acidental; ela reforça a ideia de que a implementação pode ser estudada em múltiplas escalas e com diferentes estratégias metodológicas, desde que o foco permaneça nos processos concretos que transformam decisões em ação pública. \parencite{lotta2019_teorias_analises}

Os achados e contribuições do livro devem ser lidos menos como descoberta empírica singular e mais como intervenção disciplinar. A coletânea mostra que a implementação pode iluminar temas frequentemente apartados na análise de políticas: capacidades burocráticas, coordenação interfederativa, mudanças de diretrizes ao longo da execução, interseções entre implementação e avaliação, produção de vulnerabilidades e consequências não previstas. Ao reunir esses eixos em um mesmo volume, o livro contribui para institucionalizar no Brasil um vocabulário analítico compartilhado e para demonstrar que a implementação não é apenas uma subárea especializada, mas um ponto de passagem necessário para compreender desempenho estatal, desigualdade de acesso e variação territorial das políticas. Sua contribuição para a disciplina reside precisamente em deslocar a análise do Estado abstrato para o Estado em operação. \parencite{lotta2019_teorias_analises}

Há também implicações substantivas para a administração pública e para os estudos de políticas públicas. Ao insistir que a política se transforma na implementação, o livro relativiza a suficiência de reformas centradas exclusivamente em desenho institucional formal, metas ou normativas. Ele sugere que resultados dependem da combinação entre regras, capacidades, relações interorganizacionais, perfis burocráticos, contextos territoriais e formas de interação com usuários. Essa inflexão tem efeito analítico e normativo: analiticamente, impede leituras simplistas sobre sucesso e fracasso de políticas; normativamente, aponta para a necessidade de pensar governança, coordenação, formação burocrática e mecanismos de aprendizagem a partir do terreno efetivo da execução, e não apenas de seus planos ideais. \parencite{lotta2019_teorias_analises}

Seus limites e tensões decorrem justamente de sua ambição de mapear e consolidar um campo. Como obra coletiva e programática, o livro privilegia amplitude sobre unificação teórica forte. Isso é produtivo para abrir agenda, mas deixa em segundo plano a resolução de tensões entre abordagens distintas de implementação, entre escalas analíticas e entre diferentes concepções de causalidade. Além disso, o enquadramento pelo ciclo de políticas, ainda que assumido como heurístico, mantém certa segmentação entre formulação, implementação e avaliação que alguns capítulos procuram problematizar. Em outros termos, a coletânea avança ao mostrar a centralidade da implementação, mas também evidencia que o próprio campo ainda debate como articular essa centralidade com uma visão menos faseada e mais relacional da política pública. \parencite{lotta2019_teorias_analises}

No diálogo com o restante do corpus, esta obra funciona como peça de enquadramento analítico. O estudo sobre os sistemas nacionais de políticas públicas e o Suas aprofunda, em chave institucional e federativa, questões apenas abertas na coletânea sobre coordenação, arranjos e efeitos da estruturação nacional da implementação; pode ser lido, portanto, como um desdobramento empírico de preocupações sistematizadas no volume. \parencite{bichir2020_sistemas_nacionais} Já a análise sobre categorização de usuários radicaliza a centralidade da burocracia de nível de rua e da discricionariedade, mostrando como práticas cotidianas de implementação produzem diferenças sociais concretas; ela opera em escala mais micro e relacional, mas confirma a aposta da coletânea de que a implementação é lugar de produção da política, e não simples sua aplicação. \parencite{lotta2020_categorizando_usuarios} Por sua vez, o texto sobre interações entre Estado e organizações religiosas desloca a agenda para formas híbridas de provisão e cooperação, expandindo o repertório da implementação para além da burocracia estatal estrita. Em conjunto, esses diálogos mostram que a coletânea não fornece uma unidade temática fechada ao corpus, mas um vocabulário comum suficientemente flexível para tornar comparáveis objetos distintos sem apagar suas diferenças substantivas. \parencite{lotta2023_jesus_speaks}



\subsection{Categorizando Usuários “Fáceis” e “Difíceis”: Práticas Cotidianas de Implementação de Políticas Públicas e a Produção de Diferenças Sociais}
\label{sec:orgc7aa3de}

O artigo parte de uma questão decisiva para o campo da implementação: como interações rotineiras entre agentes públicos e usuários, aparentemente ordinárias, podem produzir diferenciações socialmente relevantes no acesso, no atendimento e no usufruto de serviços públicos? Em vez de tratar a implementação como simples execução de normas, o texto desloca o foco para o encontro cotidiano entre Estado e população e sustenta que é nesse plano microssocial que se consolidam efeitos distributivos e simbólicos profundos. A preocupação central não é apenas mostrar que burocratas de nível de rua exercem discricionariedade, algo já amplamente reconhecido, mas examinar como esse exercício se conecta à produção institucional de diferenças sociais em contextos de forte desigualdade. O problema, portanto, não é só administrativo; é simultaneamente político, moral e relacional, pois envolve os modos pelos quais o Estado reconhece, seleciona, aproxima ou distancia sujeitos concretos em suas trajetórias de acesso a direitos. \parencite{lotta2020_categorizando_usuarios}

O objetivo e o escopo do estudo são formulados com precisão: compreender as práticas cotidianas adotadas por agentes envolvidos na implementação, em interação direta com os públicos-alvo, e analisar de que forma essas práticas promovem diferenciação social entre os atendidos. Para isso, o artigo combina dois movimentos articulados. De um lado, reconstrói e organiza vertentes da literatura internacional sobre implementação centrada nos agentes, identificando nelas recursos analíticos para pensar a relação entre práticas de rua e produção de diferenças. De outro, coloca esse debate em diálogo com um estudo empírico sobre a Estratégia Saúde da Família em São Paulo, com foco nos Agentes Comunitários de Saúde. O escopo, assim, é ao mesmo tempo teórico e empírico: não se limita a descrever um caso, nem se contenta com uma revisão conceitual abstrata, mas usa o caso para testar e refinar um enquadramento analítico sobre discricionariedade, categorização e desigualdade. \parencite{lotta2020_categorizando_usuarios}

A tese principal do artigo é que as diferenças sociais penetram o mundo das políticas públicas por meio de práticas de categorização e julgamento realizadas pelos agentes de implementação. Esses agentes constroem segmentações informais do público, distinguindo usuários “fáceis” e “difíceis”, e essa distinção orienta o modo como o trabalho é organizado e como os serviços são efetivamente entregues. O ponto decisivo é que tal categorização não deriva apenas de critérios técnicos ou legais, mas de uma combinação entre imperativos funcionais do trabalho e avaliações morais sobre os usuários. Assim, a implementação deixa de ser entendida apenas como arena de adaptação local de regras e passa a ser vista como espaço de produção de fronteiras simbólicas que podem reforçar vulnerabilidades preexistentes. O argumento é forte porque desloca o debate da mera existência da discricionariedade para a forma concreta pela qual ela classifica, hierarquiza e distribui oportunidades de acesso. \parencite{lotta2020_categorizando_usuarios}

No plano conceitual, o texto mobiliza um conjunto particularmente fecundo de categorias. A primeira é a discricionariedade dos burocratas de nível de rua, retomada como capacidade prática de adaptar regras, prioridades e rotinas diante de condições de trabalho marcadas por complexidade, escassez e incerteza. A segunda é o repertório analítico associado a coping, creaming e gatekeeping, usado para mostrar que a seleção informal de casos, a priorização de situações mais manejáveis e a imposição de barreiras de acesso podem ser burocraticamente funcionais, ainda que distributivamente problemáticas. A terceira, que constitui o aporte mais original do texto, é a noção de fronteiras simbólicas, empregada para iluminar como classificações elaboradas no cotidiano estatal não apenas organizam o trabalho, mas também produzem distinções socialmente significativas entre tipos de usuários. Com isso, o artigo articula literatura de implementação e sociologia da diferenciação social sem reduzir uma à outra. \parencite{lotta2020_categorizando_usuarios,lotta2019_teorias_analises}

Essa articulação conceitual também permite ao artigo propor uma distinção relevante entre diferença social e desigualdade social. O texto não trata toda classificação como desigualdade já consumada, mas como parte de um processo em que fronteiras simbólicas podem se converter em fronteiras sociais mais duras quando passam a orientar atendimento, atenção, tempo de espera, esforço institucional e qualidade da relação com o usuário. Esse cuidado analítico é importante porque evita duas simplificações: a de supor que toda discricionariedade é necessariamente arbitrária, e a de imaginar que a implementação reproduz desigualdades apenas por defeitos normativos ou falhas de desenho. O artigo mostra, ao contrário, que parte importante do problema emerge na tradução cotidiana das políticas, onde vulnerabilidades dos usuários e constrangimentos do trabalho se encontram e ganham forma classificatória. \parencite{lotta2020_categorizando_usuarios}

Em termos de método, o estudo adota uma estratégia qualitativa centrada nas narrativas dos Agentes Comunitários de Saúde sobre suas rotinas, suas percepções dos usuários e as decisões práticas que tomam no cotidiano da implementação. O foco nas narrativas é coerente com a ambição do artigo, porque o que está em jogo não é apenas registrar comportamentos observáveis, mas compreender os esquemas classificatórios que orientam a ação. Ao analisar como os agentes descrevem usuários, justificam prioridades e atribuem sentido às dificuldades do trabalho, os autores conseguem captar a dimensão cognitiva e moral da implementação. O caso da Estratégia Saúde da Família em São Paulo é particularmente adequado para isso, pois se trata de uma política intensiva em contato interpessoal, voltada a públicos heterogêneos e frequentemente marcados por vulnerabilidades múltiplas, o que torna visíveis os dilemas de triagem, priorização e vínculo. \parencite{lotta2020_categorizando_usuarios}

Os achados mais relevantes indicam que a distinção entre usuários fáceis e difíceis opera como uma gramática prática para organizar o trabalho sob pressão. Usuários considerados fáceis são, em geral, aqueles percebidos como responsivos, aderentes, compreensíveis e cooperativos, isto é, aqueles com os quais a ação estatal flui com menos atrito. Já os difíceis são associados a resistência, baixa adesão, comportamentos tidos como inadequados, maior imprevisibilidade ou necessidade de esforço adicional. O ponto central é que essas classificações não são neutras: elas tendem a influenciar a intensidade da atenção, a disposição para insistir no acompanhamento e a própria qualidade do vínculo entre agente e usuário. O que se apresenta como solução pragmática para fazer o trabalho caber nas condições existentes pode converter-se em mecanismo de seletividade cotidiana. \parencite{lotta2020_categorizando_usuarios}

Uma contribuição analítica importante do texto está em mostrar que essa seletividade não decorre exclusivamente de preconceito explícito nem exclusivamente de cálculo instrumental. Ela se forma na intersecção entre elementos funcionais e elementos morais. Funcionais, porque os agentes precisam gerir escassez, volume de demandas e metas de factibilidade do trabalho. Morais, porque avaliam comportamentos, disposições e estilos de relação dos usuários, distinguindo aqueles vistos como merecedores, comprometidos ou disciplinados daqueles percebidos como refratários, desorganizados ou pouco colaborativos. Ao explicitar essa dupla dimensão, o artigo refina a análise da discricionariedade e ajuda a compreender por que mecanismos de diferenciação persistem mesmo em políticas orientadas por direitos e universalismo. O problema não está apenas no desenho da política, mas na coprodução cotidiana de critérios implícitos de merecimento e manejabilidade. \parencite{lotta2020_categorizando_usuarios,lotta2019_teorias_analises}

Para a disciplina, o artigo oferece uma contribuição substantiva ao recolocar a implementação como espaço de produção de desigualdades, e não apenas como etapa posterior ao desenho ou à decisão. Ele sugere que os efeitos das políticas dependem, em grande medida, das formas concretas de encontro entre agentes e cidadãos, onde se definem graus de acolhimento, insistência, tradução, flexibilização e reconhecimento. Essa inflexão tem implicações importantes para estudos de políticas públicas no Brasil, especialmente porque conecta a literatura sobre burocracias de nível de rua a contextos de desigualdade estrutural. Em vez de assumir um usuário abstrato e homogêneo, o texto mostra que a heterogeneidade social entra no aparato estatal por meio de classificações práticas. Com isso, a implementação aparece como mecanismo de reprodução ou mitigação de desvantagens, dependendo de como essas classificações são produzidas e institucionalmente contidas. \parencite{lotta2020_categorizando_usuarios}

Os limites e tensões do artigo também são relevantes. O foco nas narrativas dos agentes, embora metodologicamente produtivo, privilegia a perspectiva de quem implementa e deixa menos desenvolvida a experiência dos usuários classificados e os efeitos observáveis dessas classificações sobre trajetórias concretas de acesso e cuidado. Além disso, a ênfase na construção de categorias informais ilumina bem o plano microinteracional, mas abre espaço para investigar de forma mais sistemática como organizações, metas, supervisão e arranjos institucionais mais amplos estabilizam, inibem ou transformam tais práticas. Há ainda uma tensão constitutiva no argumento: ao mostrar que certas classificações são funcionalmente racionais para trabalhadores submetidos a escassez, o texto evita moralizar os agentes, mas também revela a dificuldade de distinguir onde termina a adaptação prática e onde começa a cristalização de injustiças cotidianas. Essa tensão não enfraquece o artigo; ao contrário, indica a complexidade do fenômeno que ele torna visível. \parencite{lotta2020_categorizando_usuarios}

No diálogo com o restante do corpus, o artigo se aproxima do livro organizado por Gabriela Lotta ao compartilhar a premissa de que a implementação precisa ser compreendida a partir das práticas, interações e margens de decisão dos agentes, e não apenas por desenhos normativos ou modelos lineares de execução. Mas aqui há um avanço específico: a coletânea oferece um quadro mais amplo das teorias e análises de implementação, ao passo que este texto aprofunda uma dimensão particular, a produção de diferenças sociais por meio da categorização cotidiana dos usuários. Em relação ao estudo sobre sistemas nacionais de políticas públicas, a diferença de escala é decisiva. Enquanto a análise sobre o Suas enfatiza coordenação federativa, capacidades estatais e efeitos institucionais de arranjos nacionais, este artigo mostra que mesmo políticas fortemente institucionalizadas dependem de mediações locais e classificações de rua que modulam seus efeitos reais. Já o texto sobre interações entre Estado e organizações religiosas opera em outro terreno, o das relações interorganizacionais e da hibridização; ainda assim, há uma afinidade de fundo: ambos revelam que a implementação é atravessada por julgamentos, encaixes práticos e formas de ajuste situacional que não podem ser deduzidos apenas do desenho formal. O diálogo entre os textos, portanto, não produz unidade artificial, mas evidencia que a implementação é multiescalar: institucional, organizacional e interacional ao mesmo tempo. \parencite{lotta2020_categorizando_usuarios,lotta2019_teorias_analises,bichir2020_sistemas_nacionais,lotta2023_jesus_speaks}



\subsection{“Jesus Speaks Better”: Interactions Between State and Faith-Based Organizations in Service Implementation}
\label{sec:orge7ad175}

O artigo toma como problema central a cooperação entre organizações estatais e organizações confessionais na implementação de serviços públicos, deslocando o foco de debates normativos abstratos sobre secularismo, privatização ou terceirização para a dinâmica concreta de interação entre atores que operam com lógicas institucionais distintas. O caso escolhido — comunidades terapêuticas religiosas contratadas pelo governo paulista para atendimento de usuários problemáticos de álcool e outras drogas — permite aos autores examinar uma zona de fricção particularmente densa, porque nela se cruzam moralidades religiosas, burocracias públicas, regimes de financiamento, expectativas sobre cuidado e concepções divergentes sobre tratamento. A questão não é apenas se tais arranjos funcionam, mas como se tornam possíveis, por quais mediações se estabilizam e que efeitos produzem sobre as organizações e sobre o próprio serviço implementado.\parencite{lotta2023_jesus_speaks}

O objetivo e o escopo do texto são formulados com precisão: identificar os fatores que permitem a cooperação entre Estado e organizações confessionais na implementação de serviços e analisar as consequências dessa cooperação para as próprias organizações religiosas, para as organizações estatais e para os serviços ofertados. Em vez de tratar as entidades religiosas como meros prestadores contratados, o artigo investiga a coprodução relacional do arranjo implementador. Com isso, delimita um escopo intermediário: não pretende oferecer uma teoria geral sobre religião e políticas públicas, nem uma avaliação de impacto dos resultados terapêuticos, mas reconstruir os mecanismos institucionais e interacionais que sustentam a cooperação em contexto de forte heterogeneidade valorativa. Essa delimitação é importante porque preserva o caráter processual do argumento e evita reduzir o fenômeno a um juízo binário sobre êxito ou fracasso.\parencite{lotta2023_jesus_speaks}

A tese principal sustenta que a cooperação entre Estado e organizações confessionais depende de formas específicas de ajuste institucional e de mediação que tornam compatíveis, ao menos parcialmente, universos cognitivos e valorativos distintos. O ponto decisivo do artigo é que esse ajuste não implica simples submissão de um polo ao outro. Ao contrário, as interações geram um processo de constituição mútua: o Estado incorpora categorias, sensibilidades e quadros interpretativos provenientes das entidades religiosas, enquanto estas internalizam exigências, linguagens e modos de organização próprios do aparato estatal. O título do artigo sintetiza bem essa tese, ao indicar que determinadas categorias religiosas podem adquirir eficácia prática na relação com os usuários e, por isso, ganhar reconhecimento, tradução ou tolerância por parte do Estado. A cooperação, portanto, não nasce da eliminação das diferenças, mas da construção situada de encaixes entre elas.\parencite{lotta2023_jesus_speaks}

A arquitetura conceitual do texto combina três eixos. O primeiro é o de organizações confessionais, definidas não apenas pela vinculação formal a denominações religiosas, mas pela presença constitutiva de princípios religiosos em sua identidade institucional e em suas práticas. O segundo é a noção de hibridização, entendida como sobreposição de setores, mecanismos de governança e lógicas organizacionais em arranjos de provisão de serviços públicos. O terceiro, mais original no modo como é mobilizado pelos autores, articula interação socioestatal e ajuste institucional. A partir dessa chave, a análise rompe com a dicotomia rígida entre Estado e sociedade civil e passa a observar padrões relacionais, sedimentações históricas e acomodações práticas que permitem a cooperação. O conceito de ajuste institucional deixa de ser um simples atributo estrutural e aparece como resultado de interações que produzem encaixes cognitivos, normativos e organizacionais.\parencite{lotta2023_jesus_speaks}

Um mérito analítico relevante está em não tratar a hibridização como categoria celebratória. O artigo reconhece que a cooperação entre setores pode ampliar capilaridade, continuidade e capacidade de atendimento, mas também pode comprometer princípios públicos, gerar deslocamentos organizacionais e introduzir tensões em torno de universalismo, inclusão e critérios legítimos de intervenção. Assim, as comunidades terapêuticas não são apresentadas nem como solução natural para a insuficiência estatal nem como deformação intrinsecamente ilegítima do serviço público. Elas aparecem como arranjos ambivalentes, cuja inteligibilidade depende da observação dos mecanismos de mediação que permitem a convivência entre exigências burocráticas, valores religiosos e expectativas sobre eficácia do cuidado. Essa recusa da polarização simplificadora dá densidade ao argumento e o aproxima de uma sociologia da implementação atenta à prática.\parencite{lotta2023_jesus_speaks}

Do ponto de vista metodológico, o estudo baseia-se em pesquisa qualitativa centrada em um estudo de caso, apoiado em análise documental e em 67 entrevistas com representantes das organizações confessionais e gestores públicos envolvidos na implementação. A escolha é coerente com o problema formulado: se o interesse está nos fatores que viabilizam a cooperação e em seus efeitos sobre atores e serviços, é necessário captar interpretações, conflitos, justificações e rotinas que dificilmente apareceriam em desenhos exclusivamente quantitativos. A estratégia argumentativa é processual e relacional. Em vez de inferir a cooperação a partir do contrato formal, os autores reconstroem como ela se efetiva no cotidiano por meio de acomodações, traduções e disputas. O resultado é uma explicação baseada em mecanismos, mais do que em correlações agregadas.\parencite{lotta2023_jesus_speaks}

Os achados centrais indicam que a cooperação se sustenta em diferentes formas de mediação, descritas pelos autores sobretudo como mediações simbólicas e posicionais. As primeiras dizem respeito à tradução entre repertórios de sentido, permitindo que categorias religiosas e categorias estatais se tornem parcialmente inteligíveis umas para as outras. As segundas remetem à ocupação de posições, trajetórias ou lugares organizacionais capazes de conectar os dois mundos institucionais. Essas mediações produzem ajuste institucional porque reduzem atritos, criam confiança operacional e fornecem gramáticas compartilháveis para a ação. Ao mesmo tempo, o texto mostra que tais encaixes jamais são completos: eles estabilizam a cooperação, mas também redefinem as fronteiras do conflito. Em outras palavras, a mediação não elimina tensões; ela torna a cooperação possível apesar delas e, por vezes, as reconfigura.\parencite{lotta2023_jesus_speaks}

A principal contribuição do artigo para a disciplina está em mostrar que a implementação de políticas públicas não pode ser compreendida apenas como execução de desenhos normativos ou como problema administrativo de coordenação. Ela envolve disputa e circulação de valores, categorias e formas de ver os usuários, os serviços e a finalidade da intervenção. Ao analisar a copresença de Estado e religião na provisão pública, o texto amplia a agenda da implementação ao incorporar, de modo mais explícito, a dimensão cognitiva e moral dos arranjos implementadores. Isso é particularmente importante para a administração pública e para a análise de políticas porque evidencia que capacidades estatais, contratualização e governança dependem também de traduções simbólicas e de atores capazes de operar entre mundos institucionais heterogêneos.\parencite{lotta2023_jesus_speaks}Em diálogo com a agenda mais ampla de implementação sistematizada em \parencite{lotta2019_teorias_analises}

Há ainda uma contribuição substantiva para o entendimento de como categorias sobre públicos e intervenções podem ser coproduzidas nos encontros entre organizações. Embora o foco do texto não seja a classificação de usuários, a análise sugere que repertórios religiosos e burocráticos moldam a leitura sobre quem pode ser ajudado, de que modo e sob quais condições. Nesse sentido, o artigo conversa de forma fecunda com estudos que mostram como a implementação produz diferenças sociais por meio de categorização prática dos destinatários. A diferença é que, aqui, a ênfase recai menos sobre o encontro entre burocrata e usuário individual e mais sobre o encontro entre organizações dotadas de matrizes morais e institucionais distintas. O ganho analítico está em deslocar a produção de categorias do nível micro da interação de rua para o nível meso das interfaces organizacionais.Essa aproximação é visível quando se compara o artigo com \parencite{lotta2020_categorizando_usuarios} e com \parencite{lotta2023_jesus_speaks}

Entre os limites e tensões do artigo, sobressai o fato de que a força explicativa do estudo está ancorada em um caso específico, marcado por um setor de política pública particularmente sensível do ponto de vista moral e por organizações religiosas com trajetória própria no contexto brasileiro. Isso não reduz sua relevância, mas recomenda cautela na generalização para outros domínios de provisão, outras formas de organização confessional ou arranjos menos dependentes de mediação valorativa intensa. Além disso, ao privilegiar as condições de cooperação, o texto ilumina menos os efeitos distributivos e normativos dessa cooperação do ponto de vista dos usuários: permanece em segundo plano como diferentes concepções de cuidado, disciplina e recuperação repercutem sobre direitos, inclusão ou assimetrias de poder. A própria noção de ajuste institucional, muito produtiva para explicar coordenação, pode deixar menos visíveis situações em que o encaixe se produz às custas de silenciamentos ou acomodações assimétricas.\parencite{lotta2023_jesus_speaks}

No diálogo com o restante do corpus, o artigo se distingue por deslocar a análise da implementação do interior do aparato estatal para as interfaces entre Estado e sociedade civil organizada, especialmente em arranjos híbridos. Com o texto sobre teorias e análises da implementação, compartilha a recusa de modelos lineares e a defesa de uma visão processual, relacional e multinível da ação pública. Com o estudo sobre categorização de usuários, compartilha a atenção às dimensões cognitivas e morais que atravessam a implementação, embora as localize em outro plano analítico. E com a discussão sobre sistemas nacionais e seus efeitos na implementação, estabelece um contraste útil: enquanto aquele enfatiza capacidades, coordenação federativa e indução institucional em políticas nacionalmente estruturadas, este mostra que, mesmo quando há contratualização e regras, o funcionamento efetivo da implementação depende de mediações situadas entre organizações portadoras de lógicas heterogêneas. Não há unidade artificial entre os textos, mas há um eixo comum robusto: a implementação aparece, em todos, como processo de produção concreta da política por meio de interações, classificações, encaixes institucionais e disputas sobre o que conta como ação pública legítima.\parencite{lotta2019_teorias_analises}; \parencite{lotta2020_categorizando_usuarios}; \parencite{bichir2020_sistemas_nacionais}; \parencite{lotta2023_jesus_speaks}



\section{Comparação entre os textos}
\label{sec:orgbf3619f}

\begin{landscape}
\begingroup
\scriptsize
\setlength{\tabcolsep}{3pt}
\renewcommand{\arraystretch}{1.12}
\begin{longtblr}[caption={Síntese comparativa dos principais eixos transversais do corpus},label={tbl-comparacao-corpus}]{%
  width=\linewidth,
  colspec={X[1.0,l] X[1.2,l] X[1.15,l] X[1.15,l] X[1.25,l] X[1.15,l]},
  rowhead=1,
  rows={valign=t},
  row{1}={font=\bfseries},
  hline{1,Z}={0.08em},
  hline{2}={0.04em},
}
Texto & Objeto e problema & Escala e método & Conceitos centrais & Contribuição & Limite/tensão \\
Estudo sobre o Suas & Efeitos de um sistema nacional sobre a implementação municipal e a redução de desigualdades territoriais & Escala federativa e municipal; análise comparativa com indicadores nacionais de implementação e capacidade & sistemas nacionais, coordenação federativa, indução federal, capacidades institucionais, governança multinível & Mostra que arranjos nacionais podem induzir expansão e convergência da implementação sem eliminar a relevância das capacidades locais & Capta melhor estrutura e oferta do que interações concretas, qualidade do atendimento e experiência dos usuários \\
Coletânea sobre implementação no Brasil & Consolidação do campo de estudos da implementação e demonstração de sua centralidade analítica & Escala teórico-programática e multissetorial; síntese conceitual e pluralidade de estudos empíricos & implementação, burocratas de nível de rua, discricionariedade, capacidades estatais, coordenação, arranjos institucionais & Oferece vocabulário comum para compreender a política pública em operação, para além do desenho formal & Amplitude maior que integração teórica; abre agendas mais do que resolve tensões entre abordagens \\
Artigo sobre categorização de usuários & Como práticas cotidianas de agentes produzem diferenças sociais no atendimento & Escala microinteracional; estudo qualitativo com foco nas narrativas de implementadores & discricionariedade, categorização, usuários fáceis e difíceis, fronteiras simbólicas, coping, creaming, gatekeeping & Explicita como a implementação pode reproduzir ou aprofundar desigualdades por classificações morais e funcionais & Privilegia a perspectiva dos agentes e menos os efeitos observados nas trajetórias dos usuários \\
Artigo sobre Estado e organizações religiosas & Condições de cooperação entre Estado e organizações confessionais na implementação de serviços & Escala interorganizacional; estudo de caso qualitativo com documentos e entrevistas & organizações confessionais, hibridização, mediações simbólicas, mediações posicionais, ajuste institucional & Amplia o campo da implementação para interfaces híbridas entre Estado e sociedade civil e para a dimensão moral dos arranjos & Generalização limitada e menor desenvolvimento dos efeitos normativos e distributivos para os usuários \\
\end{longtblr}
\endgroup
\end{landscape}

Tomados em conjunto, os quatro textos convergem em um ponto decisivo: a implementação não é tratada como fase residual, mecânica ou subordinada da política pública, mas como espaço de produção efetiva da ação estatal, no qual normas, organizações, burocracias, valores e interações reconfiguram o que a política é em termos concretos. Essa convergência, contudo, não deve ser confundida com homogeneidade temática ou teórica plena. O corpus é internamente heterogêneo e sua força reside justamente em mostrar que a implementação pode ser observada em escalas muito distintas — sistêmica, programática, interacional e interorganizacional — sem que uma delas esgote o fenômeno. A coletânea organizada por Gabriela Lotta fornece o enquadramento mais geral dessa premissa, ao sustentar que a política formulada não coincide com a política efetivamente produzida \parencite{lotta2019_teorias_analises}. Os demais trabalhos, por sua vez, demonstram empiricamente essa proposição em registros analíticos diversos: um pela via da coordenação federativa e dos efeitos de sistemas nacionais \parencite{bichir2020_sistemas_nacionais}; outro pela via da categorização cotidiana de usuários e da produção de diferenças sociais \parencite{lotta2020_categorizando_usuarios}; e outro ainda pelas mediações que tornam possível a cooperação entre Estado e organizações confessionais \parencite{lotta2023_jesus_speaks}.

A primeira diferença transversal substantiva diz respeito ao objeto de análise. O estudo sobre o Suas examina um arranjo institucional nacionalmente estruturado e pergunta pelos efeitos desse desenho sobre a implementação municipal, isto é, pelo modo como regras federais, transferências, pactuações e parâmetros nacionais induzem expansão e convergência na provisão local \parencite{bichir2020_sistemas_nacionais}. Já a coletânea sobre implementação não se organiza em torno de um único objeto substantivo, mas de um objeto disciplinar: o próprio campo dos estudos de implementação, sua autonomia relativa e sua capacidade de iluminar o funcionamento concreto das políticas públicas \parencite{lotta2019_teorias_analises}. O artigo sobre categorização de usuários desloca a lente para o encontro entre agentes e cidadãos, focalizando o trabalho cotidiano de classificação que transforma heterogeneidade social em tratamento burocrático diferencial \parencite{lotta2020_categorizando_usuarios}. Por fim, o texto sobre organizações religiosas examina as interfaces entre Estado e sociedade civil organizada, analisando como a implementação se estabiliza em arranjos híbridos atravessados por lógicas normativas distintas \parencite{lotta2023_jesus_speaks}. Em termos comparativos, isso significa que o corpus não trata de um único problema empírico, mas de múltiplas entradas para uma mesma questão mais ampla: como a ação pública ganha forma na prática.

Essa pluralidade de objetos corresponde a diferenças igualmente importantes de escala analítica. O artigo sobre o Suas opera em escala macro e meso, porque relaciona a arquitetura federativa de um sistema nacional a padrões agregados de implementação municipal \parencite{bichir2020_sistemas_nacionais}. O livro organizado por Lotta transita entre escalas, mas sua função principal é meta-analítica: oferecer categorias que permitam conectar burocracias, arranjos, coordenação, capacidades e efeitos não previstos em diferentes contextos \parencite{lotta2019_teorias_analises}. O estudo sobre usuários fáceis e difíceis se concentra em escala micro, nas práticas situadas e nos esquemas classificatórios que organizam o trabalho dos implementadores \parencite{lotta2020_categorizando_usuarios}. O artigo sobre organizações confessionais trabalha em escala meso, privilegiando relações interorganizacionais e mecanismos de mediação entre diferentes matrizes institucionais \parencite{lotta2023_jesus_speaks}. A comparação mostra, portanto, que o corpus é menos um conjunto de respostas concorrentes do que uma composição multiescalar. Se lidos em sequência, os textos sugerem que a implementação precisa ser pensada como fenômeno simultaneamente estruturado por instituições amplas, mediado por organizações e produzido em interações concretas.

Os métodos empregados reforçam essa heterogeneidade. O estudo do Suas é o que mais se aproxima de uma estratégia comparativa nacional baseada em indicadores agregados, buscando medir variação e convergência territorial ao longo do tempo \parencite{bichir2020_sistemas_nacionais}. A coletânea, por definição, não adota um método único; ela combina síntese teórica com demonstrações empíricas plurais, o que a torna mais útil como mapa de abordagens do que como intervenção causal unificada \parencite{lotta2019_teorias_analises}. Os artigos sobre categorização de usuários e sobre organizações religiosas recorrem a estratégias qualitativas, mas com focos distintos: no primeiro, as narrativas de agentes permitem reconstruir esquemas cognitivos e morais de classificação \parencite{lotta2020_categorizando_usuarios}; no segundo, entrevistas e documentos sustentam uma análise processual de mediações e ajustes institucionais \parencite{lotta2023_jesus_speaks}. Não há aqui superioridade intrínseca de um método sobre outro. O ponto comparativo relevante é que cada estratégia responde a um tipo diferente de pergunta. O desenho agregativo permite captar padrões de expansão e convergência; os desenhos qualitativos permitem captar mecanismos, classificações, conflitos e traduções invisíveis em indicadores nacionais. A articulação possível entre os textos depende de reconhecer essa divisão de trabalho analítica, e não de exigir de cada um aquilo que apenas outro poderia oferecer.

No plano conceitual, há convergências claras, mas também arquiteturas distintas. A coletânea oferece o repertório mais abrangente, com categorias como implementação, discricionariedade, burocracias de nível de rua, burocracias de médio escalão, capacidades estatais, coordenação e arranjos institucionais \parencite{lotta2019_teorias_analises}. O texto sobre o Suas apropria-se de parte desse vocabulário, mas o reorganiza em torno de sistemas nacionais de políticas públicas, coordenação federativa, indução federal e governança multinível \parencite{bichir2020_sistemas_nacionais}. O artigo sobre categorização de usuários radicaliza o foco na discricionariedade, articulando-o a categorias mais finas, como coping, creaming, gatekeeping e fronteiras simbólicas, para mostrar como o tratamento burocrático diferencia sujeitos \parencite{lotta2020_categorizando_usuarios}. Já o estudo sobre organizações religiosas desloca o vocabulário para hibridização, mediações simbólicas, mediações posicionais e ajuste institucional, expandindo o campo da implementação para relações entre organizações portadoras de racionalidades distintas \parencite{lotta2023_jesus_speaks}. Comparativamente, isso significa que o corpus não compartilha uma gramática única, mas um núcleo comum: todos os textos rejeitam a ideia de que regras bastam para explicar a política em ação. Divergem, porém, quanto ao local privilegiado onde essa insuficiência se manifesta: no federalismo, nas burocracias, nas interações com usuários ou nas interfaces entre organizações.

Os pressupostos teóricos também se aproximam sem se confundirem. Em todo o corpus há uma crítica, explícita ou implícita, às leituras lineares do ciclo de políticas quando ele é tomado como descrição empírica simples. A coletânea formula isso de modo mais programático ao insistir que implementação não é mera execução \parencite{lotta2019_teorias_analises}. O artigo sobre o Suas, embora menos centrado em discussão epistemológica, compartilha um institucionalismo orientado por efeitos de arranjos e capacidades, no qual a política é resultado de interação entre regras, incentivos e condições locais \parencite{bichir2020_sistemas_nacionais}. O texto sobre categorização de usuários opera com pressupostos mais próximos de uma sociologia interacional e moral da ação pública, mostrando que categorias e julgamentos cotidianos são constitutivos da distribuição concreta de direitos \parencite{lotta2020_categorizando_usuarios}. O artigo sobre organizações religiosas, por sua vez, enfatiza a coprodução relacional e os processos de tradução entre lógicas institucionais heterogêneas, recusando tanto a separação rígida entre Estado e sociedade civil quanto a ideia de cooperação como simples alinhamento funcional \parencite{lotta2023_jesus_speaks}. A principal divergência teórica, assim, não está em um desacordo frontal, mas na ênfase causal privilegiada por cada texto: estruturas e incentivos, vocabulário de campo, classificação prática ou mediação institucional.

Quanto às convergências empíricas reais, duas merecem destaque. A primeira é que todos os textos mostram que a implementação é atravessada por mediações. No estudo do Suas, a mediação aparece na relação entre coordenação federal e capacidade municipal \parencite{bichir2020_sistemas_nacionais}. No livro, ela aparece na multiplicidade de atores, escalas e arranjos que tornam a política diferente de seu desenho \parencite{lotta2019_teorias_analises}. No artigo sobre usuários, a mediação assume a forma de julgamentos e classificações produzidos por agentes de linha de frente \parencite{lotta2020_categorizando_usuarios}. No texto sobre organizações confessionais, surge como tradução entre repertórios estatais e religiosos \parencite{lotta2023_jesus_speaks}. A segunda convergência é a recusa de imaginar o Estado como ator unitário. Em todos os casos, o Estado aparece distribuído em níveis federativos, burocracias, organizações, contratos e interações. Isso constitui contribuição importante para a disciplina, porque desloca a análise de modelos jurídicos ou formalistas para configurações efetivas de ação pública.

As divergências reais, por sua vez, dizem respeito sobretudo ao que cada texto considera decisivo para explicar variação na implementação. No artigo sobre o Suas, o ponto crucial está na capacidade dos sistemas nacionais de induzir comportamentos e reduzir desigualdades territoriais, ainda que sem anular diferenças de capacidade local \parencite{bichir2020_sistemas_nacionais}. No artigo sobre categorização de usuários, o núcleo explicativo desloca-se para o nível microssocial, mostrando que políticas formalmente orientadas por direitos podem ser moduladas por classificações morais e funcionais do público \parencite{lotta2020_categorizando_usuarios}. Já o estudo sobre organizações religiosas sugere que a eficácia ou estabilidade da implementação depende, em certos contextos, menos da aderência a regras uniformes e mais da existência de encaixes relacionais entre lógicas institucionais distintas \parencite{lotta2023_jesus_speaks}. A coletânea não arbitra essa divergência; ao contrário, legitima sua coexistência ao apresentar a implementação como campo plural \parencite{lotta2019_teorias_analises}. Comparativamente, isso impede qualquer conclusão simplista do tipo “a implementação se explica por X”. O corpus sugere, antes, que o peso relativo de instituições, capacidades, discricionariedade e mediações varia conforme o objeto e a escala.

Em termos de contribuição para a disciplina, cada texto ocupa um lugar específico e complementar. A coletânea tem função de consolidação e ensino do campo, ao sistematizar categorias e demonstrar que implementação é dimensão constitutiva da análise de políticas públicas no Brasil \parencite{lotta2019_teorias_analises}. O artigo sobre o Suas contribui ao conectar estudos de implementação com federalismo e coordenação de políticas sociais, oferecendo evidência de que sistemas nacionais importam para padrões territoriais de provisão \parencite{bichir2020_sistemas_nacionais}. O estudo sobre usuários aprofunda o debate sobre burocracias de nível de rua ao mostrar que a discricionariedade não afeta apenas eficiência ou adaptação, mas também a produção institucional de diferenças sociais \parencite{lotta2020_categorizando_usuarios}. O texto sobre organizações religiosas amplia o perímetro do campo ao incorporar arranjos híbridos, coprodução interorganizacional e disputas morais na provisão de serviços \parencite{lotta2023_jesus_speaks}. Vistos em conjunto, os textos ajudam a deslocar a disciplina de uma visão centrada exclusivamente no desenho formal para uma concepção da implementação como arena de coordenação, interpretação, classificação e tradução.

Os limites comparativos do corpus também precisam ser explicitados para evitar unidade artificial. Em primeiro lugar, trata-se de textos com naturezas diferentes: uma coletânea programática, um estudo quantitativo-comparativo de política nacional e dois artigos qualitativos focalizados em mecanismos distintos. Isso torna inadequado exigir a mesma forma de prova, o mesmo tipo de generalização ou o mesmo nível de abstração de todos eles. Em segundo lugar, os setores e contextos empíricos não são idênticos: assistência social, saúde e comunidades terapêuticas envolvem públicos, regimes de provisão e conflitos normativos diferentes. Em terceiro lugar, nem todos os textos têm como foco principal os efeitos distributivos para os usuários, ainda que todos permitam inferências sobre eles em alguma medida. O estudo do Suas fala mais diretamente de convergência institucional e territorial do que de experiência concreta do usuário \parencite{bichir2020_sistemas_nacionais}; o artigo sobre categorização faz o movimento inverso \parencite{lotta2020_categorizando_usuarios}; e o texto sobre organizações religiosas privilegia as condições de cooperação e hibridização mais do que a avaliação dos resultados substantivos do cuidado \parencite{lotta2023_jesus_speaks}. Esses limites não enfraquecem a comparação; apenas definem seu alcance legítimo.

As articulações mais promissoras entre os textos surgem quando se aceita essa heterogeneidade como vantagem analítica. O estudo do Suas pode ser lido como demonstração de que arranjos nacionais moldam oportunidades e constrangimentos da implementação local, mas ele ganha profundidade quando colocado ao lado do artigo sobre categorização, que mostra como a política assim estruturada ainda depende de encontros concretos e julgamentos práticos \parencite{bichir2020_sistemas_nacionais,lotta2020_categorizando_usuarios}. A coletânea, por sua vez, funciona como ponte entre esses níveis ao oferecer linguagem para tratar, num mesmo horizonte, capacidades, coordenação e discricionariedade \parencite{lotta2019_teorias_analises}. Já o texto sobre organizações religiosas acrescenta uma camada que os outros apenas tangenciam: a implementação pode ocorrer em ecossistemas híbridos, nos quais nem o Estado nem a burocracia pública possuem monopólio da ação implementadora \parencite{lotta2023_jesus_speaks}. Isso sugere uma leitura transversal forte, porém não artificial: a implementação é multiescalar e relacional, sendo produzida por sistemas, organizações, burocratas e interfaces institucionais. O que varia de texto para texto é o ponto de entrada analítico escolhido para tornar essa complexidade visível.

Em síntese comparativa, o corpus não autoriza uma narrativa unificada simples, mas sustenta uma conclusão robusta para a disciplina: compreender políticas públicas exige estudar a distância produtiva entre desenho e operação. Essa distância não é mero ruído. Ela é o lugar onde se decidem expansão territorial, coordenação federativa, seletividade cotidiana, cooperação interorganizacional e, em última instância, o conteúdo efetivo dos direitos e serviços públicos. O mérito do conjunto está justamente em mostrar que essa produção da política em ato pode ser observada sob diferentes lentes sem perder densidade analítica. Em vez de forçar convergência temática total, a comparação mais fiel reconhece que os textos se complementam por justaposição crítica: um ilumina os efeitos estruturais dos sistemas nacionais \parencite{bichir2020_sistemas_nacionais}, outro organiza o vocabulário e a agenda do campo \parencite{lotta2019_teorias_analises}, outro revela a microprodução de diferenças sociais \parencite{lotta2020_categorizando_usuarios}, e outro expande o debate para arranjos híbridos e moralmente densos de implementação \parencite{lotta2023_jesus_speaks}. A articulação possível entre eles, portanto, não é a de uma unidade temática artificial, mas a de uma cartografia sofisticada dos lugares onde a política pública efetivamente acontece.



\section{Síntese analítica}
\label{sec:org6a1d584}

Uma leitura articulada do corpus sugere que sua contribuição mais forte não está em responder a uma única pergunta substantiva, mas em desmontar uma imagem empobrecida do Estado. Nos quatro textos, o Estado não aparece como centro coerente que define e depois simplesmente executa políticas; ele surge como composição de níveis federativos, burocracias, organizações, interações e traduções. Essa mudança de perspectiva é decisiva porque desloca a implementação do lugar de “etapa posterior” para o lugar de problema constitutivo. A coletânea sobre implementação fornece a formulação mais explícita dessa inflexão ao mostrar que a política em operação não coincide com a política formulada \parencite{lotta2019_teorias_analises}. Os demais textos podem ser lidos como demonstrações empíricas dessa tese geral, cada um incidindo sobre uma camada distinta da ação pública.

Nessa chave, o estudo sobre o Suas é particularmente importante porque mostra que a implementação não é apenas fenômeno local e disperso; ela também é estruturada por arquiteturas institucionais nacionais capazes de induzir padrões de provisão e reduzir desigualdades territoriais \parencite{bichir2020_sistemas_nacionais}. Seu argumento impede duas simplificações recorrentes. A primeira seria imaginar que descentralização equivale a autonomia plena dos municípios e, portanto, à impossibilidade de coordenação substantiva. A segunda seria supor que a regulação nacional basta por si só para homogeneizar a implementação. O texto recusa ambas ao mostrar que sistemas nacionais importam justamente porque combinam indução federal e dependência de capacidades locais. A implementação aparece, então, como efeito de uma tensão produtiva entre centralização regulatória e recepção territorial, e não como mera execução administrativa comandada de cima.

O artigo sobre categorização de usuários desloca o problema para outro plano e, ao fazê-lo, corrige um possível excesso de institucionalismo presente em leituras centradas apenas em sistemas e arranjos. Mesmo quando a política é regulada, financiada e organizada, o conteúdo concreto do atendimento continua sendo produzido em encontros situados, nos quais agentes interpretam regras, priorizam casos e classificam públicos \parencite{lotta2020_categorizando_usuarios}. A grande força desse texto é mostrar que a implementação não apenas varia; ela diferencia. E diferencia não de modo aleatório, mas por gramáticas práticas que combinam conveniência funcional e julgamento moral. Com isso, o problema da implementação deixa de ser apenas o da eficiência organizacional ou da coordenação entre níveis de governo: ele passa a incluir a produção cotidiana de fronteiras entre usuários mais facilmente incorporáveis e usuários mais difíceis de acomodar nas rotinas estatais.

Esse deslocamento é relevante porque permite interpretar o corpus a partir de uma ideia comum: a implementação é sempre um processo de seleção, mas os critérios dessa seleção variam conforme a escala observada. No estudo do Suas, selecionam-se prioridades e padrões por meio de instrumentos institucionais e financeiros \parencite{bichir2020_sistemas_nacionais}; no texto sobre usuários, selecionam-se intensidades de atenção e formas de vínculo por meio de classificações práticas \parencite{lotta2020_categorizando_usuarios}; no artigo sobre organizações confessionais, selecionam-se formas legítimas de cooperação por meio de mediações que tornam compatíveis repertórios institucionais distintos \parencite{lotta2023_jesus_speaks}; e a coletânea oferece o vocabulário geral para reconhecer que essas seleções são constitutivas da política pública, não distorções marginais \parencite{lotta2019_teorias_analises}.

A partir daí, torna-se possível sustentar uma interpretação mais própria do conjunto: o corpus mostra que a implementação é o nível em que o Estado administra diferenças. Essas diferenças podem ser territoriais, como no caso dos municípios com capacidades desiguais diante de um sistema nacional \parencite{bichir2020_sistemas_nacionais}; sociais, como nas classificações entre usuários percebidos como fáceis e difíceis \parencite{lotta2020_categorizando_usuarios}; ou institucionais, como nas fricções entre racionalidades estatais e religiosas na coprodução de serviços \parencite{lotta2023_jesus_speaks}. Em todos os casos, a política pública só se estabiliza porque algum mecanismo de tratamento dessas diferenças entra em funcionamento. Isso significa que a implementação não é apenas mediação entre decisão e resultado; ela é o próprio trabalho de tornar governável a heterogeneidade que o desenho normativo não consegue absorver integralmente.

Essa leitura também evidencia uma tensão importante no corpus. Se, por um lado, os textos convergem ao mostrar que a implementação depende de adaptações, mediações e julgamentos, por outro eles não atribuem o mesmo valor normativo a essas mediações. No estudo do Suas, a adaptação local aparece sobretudo como condição para que a coordenação nacional se converta em expansão concreta da política \parencite{bichir2020_sistemas_nacionais}. No artigo sobre usuários, a adaptação pode ser precisamente o veículo pelo qual desigualdades são reproduzidas no atendimento \parencite{lotta2020_categorizando_usuarios}. No texto sobre organizações confessionais, a mediação é ambígua: viabiliza cooperação e capilaridade, mas também reabre o debate sobre quais valores passam a compor o serviço público \parencite{lotta2023_jesus_speaks}. A coletânea, por sua vez, legitima essa ambivalência ao apresentar a implementação como terreno de criatividade institucional e, ao mesmo tempo, de consequências não previstas \parencite{lotta2019_teorias_analises}.

Daí decorre um segundo argumento interpretativo: o corpus não autoriza uma defesa abstrata da discricionariedade, da coordenação ou da cooperação híbrida; ele exige perguntar sob quais condições esses mecanismos ampliam capacidades públicas e sob quais condições introduzem seletividades problemáticas. A coordenação federativa pode reduzir desigualdades territoriais, mas não suprime a importância da base organizacional local \parencite{bichir2020_sistemas_nacionais}. A discricionariedade pode ser inevitável para lidar com casos complexos, mas também pode converter vulnerabilidade social em baixa prioridade burocrática \parencite{lotta2020_categorizando_usuarios}. A cooperação com organizações confessionais pode ampliar adesão e inteligibilidade junto a determinados públicos, mas isso ocorre ao preço de processos de hibridização que redefinem as fronteiras do que se entende por ação pública \parencite{lotta2023_jesus_speaks}. O mérito do corpus é mostrar que esses dilemas não são exceções; são parte da normalidade da implementação.

Outra implicação importante é metodológica. Lidos conjuntamente, os textos sugerem que nenhuma escala, isoladamente, esgota a análise da implementação. Estudos agregados permitem captar padrões de expansão, convergência e efeito institucional em âmbito nacional \parencite{bichir2020_sistemas_nacionais}; análises teóricas e de balanço de campo organizam os conceitos necessários para tornar comparáveis fenômenos diversos \parencite{lotta2019_teorias_analises}; pesquisas qualitativas de nível micro revelam julgamentos e fronteiras invisíveis em indicadores padronizados \parencite{lotta2020_categorizando_usuarios}; e estudos de interfaces organizacionais mostram como valores e categorias circulam entre instituições \parencite{lotta2023_jesus_speaks}. Em vez de concorrerem, essas abordagens se completam. A síntese mais robusta do corpus depende justamente dessa composição multiescalar.

Em suma, a interpretação mais produtiva do conjunto é a de que a implementação constitui o plano em que se tornam observáveis tanto a materialidade institucional do Estado quanto sua dimensão relacional. Ela é material porque depende de sistemas, regras, financiamento, equipamentos, capacidades e arranjos organizacionais; e é relacional porque opera por classificações, traduções, vínculos, conflitos e acomodações. O corpus ganha densidade exatamente por não reduzir um desses lados ao outro. Seu ensinamento mais forte para a disciplina é que compreender políticas públicas exige acompanhar como elas são estruturadas, negociadas e redefinidas no curso de sua realização. Não há, aqui, política pública plenamente inteligível antes da implementação; há políticas que só se revelam analiticamente quando observadas no modo concreto como produzem ordem, diferença e coordenação.



\section{Considerações finais}
\label{sec:orgbf3ccc8}

As principais contribuições analíticas deste corpus decorrem de sua capacidade de requalificar a implementação como objeto central da análise de políticas públicas. Em conjunto, os textos mostram que resultados estatais não podem ser explicados apenas por desenho formal, vontade governamental ou arquitetura normativa. A implementação aparece como dimensão constitutiva da política, isto é, como espaço em que se definem alcance territorial, modos de atendimento, critérios de priorização e formas de cooperação entre atores públicos e não estatais \parencite{lotta2019_teorias_analises,bichir2020_sistemas_nacionais,lotta2020_categorizando_usuarios,lotta2023_jesus_speaks}.

Um primeiro ganho para a disciplina está na afirmação de uma perspectiva multiescalar. O estudo sobre o Suas demonstra que sistemas nacionais podem induzir expansão e convergência da implementação em contexto federativo desigual \parencite{bichir2020_sistemas_nacionais}; o artigo sobre categorização de usuários evidencia que, mesmo sob regras comuns, o atendimento é redefinido no nível das interações cotidianas \parencite{lotta2020_categorizando_usuarios}; e a análise sobre organizações confessionais mostra que a implementação também depende de encaixes entre instituições portadoras de valores distintos \parencite{lotta2023_jesus_speaks}. A coletânea organizada por Lotta oferece a moldura conceitual que torna inteligível essa variedade de escalas sem dissolver a unidade analítica do campo \parencite{lotta2019_teorias_analises}.

Um segundo ganho consiste em evidenciar que a implementação é inseparável do tratamento institucional da heterogeneidade. Os textos analisados mostram, cada um a seu modo, que a política pública se realiza ao lidar com diferenças territoriais, sociais e organizacionais. Esse resultado é especialmente relevante para a disciplina porque impede modelos excessivamente abstratos do Estado e das políticas. Em lugar de um aparato homogêneo que aplica regras a destinatários indiferenciados, o corpus apresenta um Estado que coordena, classifica, traduz e negocia. Isso aprofunda a compreensão sobre capacidades estatais, uma vez que capacidade, aqui, não significa apenas possuir recursos e normas, mas conseguir operar em contextos de diversidade e conflito.

Por fim, a principal contribuição do conjunto para a disciplina talvez seja metodológica e teórica ao mesmo tempo: ele ensina que a análise de políticas públicas ganha densidade quando acompanha a política em operação. Isso exige combinar atenção a arranjos institucionais, capacidades organizacionais, discricionariedade burocrática, relações com usuários e interfaces entre setores. A leitura do corpus, portanto, não conduz a uma teoria única da implementação, mas a um princípio analítico mais robusto: políticas públicas devem ser estudadas no ponto em que são efetivamente produzidas. É nesse plano que se tornam visíveis seus alcances, suas tensões e seus efeitos concretos sobre a organização do Estado e sobre a distribuição de direitos e serviços.

\printbibliography

\clearpage
\section{Mapa mental dos textos analisados}
\label{sec:org07c05f5}
\vspace{0.8em}
\begingroup
\centering
\includegraphics[width=\textwidth,height=0.78\textheight,keepaspectratio]{images/mapa_mental.png}
\par
\endgroup
\clearpage
\end{document}